% =============================================================================
% CASPEr@night --- Installation and First-Run Guide (Bruker / TopSpin)
% Facility-facing install guide. Everything in here is re-presented from the
% repository sources (topspin/INSTALL.md, README.md, PROTOCOL.md,
% testing/tier0_desktest.md, testing/pilot_runbook.md, uploader/, and
% topspin/spin_noise_run.py itself). Build: pdflatex x2.
% =============================================================================
\documentclass[11pt]{article}

\usepackage[T1]{fontenc}
\usepackage[utf8]{inputenc}
\usepackage{lmodern}
\usepackage{microtype}
\usepackage[margin=2.4cm]{geometry}
\usepackage{xcolor}
\usepackage{booktabs}
\usepackage{enumitem}
\usepackage{listings}
\usepackage{tcolorbox}
\usepackage[hidelinks]{hyperref}

\definecolor{cmdbg}{gray}{0.955}
\definecolor{cmdframe}{gray}{0.70}
\definecolor{boxframe}{rgb}{0.15,0.30,0.45}
\definecolor{boxbg}{rgb}{0.945,0.965,0.98}

\lstset{
  basicstyle=\ttfamily\small,
  backgroundcolor=\color{cmdbg},
  frame=single,
  rulecolor=\color{cmdframe},
  framesep=4pt,
  xleftmargin=6pt,
  xrightmargin=6pt,
  breaklines=true,
  breakatwhitespace=true,
  postbreak=\mbox{\space\space},
  columns=fullflexible,
  keepspaces=true,
  aboveskip=6pt,
  belowskip=6pt,
  literate={~}{\textasciitilde}{1}
}

\newcommand{\hnmr}{\textsuperscript{1}H}
\newcommand{\dtwo}{D\textsubscript{2}O}
\newcommand{\htwo}{H\textsubscript{2}O}
\newcommand{\ts}{TopSpin}
\newcommand{\file}[1]{\texttt{#1}}
\newcommand{\dialog}[1]{\textbf{#1}}

\setlist[itemize]{topsep=3pt, itemsep=2pt, parsep=0pt}
\setlist[enumerate]{topsep=3pt, itemsep=3pt, parsep=0pt}

\emergencystretch=1.5em

\title{\textbf{CASPEr@night --- Installation and First-Run Guide}\\
\large (Bruker / TopSpin)}
\author{John W. Blanchard (\href{mailto:jwbquantum@gmail.com}{jwbquantum@gmail.com})
\and Reza Ebadi \and Claude (Anthropic)}
\date{Software v0.3.0 --- August 2026}

\begin{document}
\maketitle

\section{What this is}

CASPEr@night is a community measurement program: any NMR facility with a
Bruker spectrometer can contribute a data point with one 5\,mm tube of water
and one \ts{} command, \texttt{xpy spin\_noise\_run}. The script records the
spin-noise feature of water --- the Gu\'eron absorption dip on room-temperature
probes, the emission bump on cryoprobes --- together with the receiver
calibrations needed to interpret it absolutely, and packs everything into a
single zip for upload. Across many facilities these bundles map the
temperature-contrast law of circuit--spin coupling over the community's probe
fleet, and they bank spin-noise-limited noise records for
fundamental-sensitivity and dark-matter (CASPEr-style) analyses.

The project page is
\href{https://night.blanchard-science.com}{night.blanchard-science.com}. The
data policy (\file{DATA\_POLICY.md} in the distribution) is short and worth the
two minutes --- your raw data remain yours, and contributing facilities are
co-authors on any network publication their data enter.

\begin{tcolorbox}[colback=boxbg, colframe=boxframe,
  title=\textbf{The five-minute version}, fonttitle=\bfseries]
\begin{enumerate}[leftmargin=1.4em]
\item Download the latest release zip:
      \url{https://github.com/CASPEr-Night/spin-noise-network/releases/latest}
\item Copy two files onto the spectrometer workstation:
      \file{topspin/spin\_noise\_run.py} $\rightarrow$
      \file{<TSHOME>/exp/stan/nmr/py/user/} and
      \file{topspin/pp/zgnoise2d} $\rightarrow$
      \file{<TSHOME>/exp/stan/nmr/lists/pp/user/}
\item Fill a 5\,mm tube with ${\sim}550\,\mu$L of plain water (any water ---
      you just have to know what it is), insert it, and open any existing
      \hnmr{} dataset in \ts.
\item Dry run first, no hardware touched:
      \texttt{xpy spin\_noise\_run simulate}
\item Real run: \texttt{xpy spin\_noise\_run} --- answer the dialogs,
      \textbf{confirm the BSMS field sweep is OFF}, walk away
      (${\sim}45$\,min default).
\item Upload from any machine with Python~3:
      \texttt{python3 uploader/upload\_bundle.py <bundle.zip>}
      (first time: copy \file{config.example.json} to \file{config.json} and
      paste the endpoint and token the coordinator sent you).
\end{enumerate}
\end{tcolorbox}

\section{Requirements}

\begin{itemize}
\item \textbf{A Bruker spectrometer running \ts{} 2.x, 3.x, or 4.x.} (\ts{} 5.0, released 2026, is untested: Bruker documents the Jython layer this software runs in as a standard component alongside the newer CPython interface, so it is expected to work --- but we have not verified it. If your console runs \ts{} 5, please tell us what happens: you would be our first.) The
      script is written for \ts's embedded Jython interpreter and avoids
      anything version-specific --- every optional command (\texttt{atma},
      \texttt{topshim}, \texttt{pulsecal}, \texttt{rga}) degrades to an
      operator dialog when missing, so old consoles work too.
\item \textbf{Any \hnmr-capable probe} --- room-temperature, nitrogen-cooled,
      or helium cryoprobe. The run asks which and records it.
\item \textbf{One 5\,mm tube of plain water}, ${\sim}550\,\mu$L. Tap,
      distilled, or \dtwo-doped all work --- \emph{as long as you know the
      composition}. The \htwo-fraction question in the dialogs is the
      measurement itself (the spin-noise feature scales with proton density),
      so if the tube contains \dtwo, know the percentage.
\item \textbf{Python 3.6 or newer on any machine that can reach the
      internet}, for the upload step only. The spectrometer host itself needs
      no internet access and no Python packages --- carry the zip on a USB
      stick if needed. Check with:
\begin{lstlisting}
python3 --version
\end{lstlisting}
\item \textbf{NTP time synchronization enabled on the \ts{} workstation}
      (it usually is by default). The run measures your console's master-clock
      offset for free by comparing the NTP-disciplined workstation clock
      against the console-timed acquisition durations --- the ``clock audit''
      in the bundle. Nothing on the spectrometer is touched or reconfigured
      --- if your workstation is air-gapped, just mention that when you
      upload.
\item \textbf{Disk and time:} the default run writes roughly 200\,MB
      (${\geq}10$\,GB free is comfortable headroom for reruns and longer noise
      blocks) and takes ${\sim}45$ minutes of magnet time, of which you are
      needed for the first ${\sim}5$ (the dialogs).
\end{itemize}

\section{Download}

Get the latest release zip from the GitHub releases page:
\begin{lstlisting}
https://github.com/CASPEr-Night/spin-noise-network/releases/latest
\end{lstlisting}
The zip is the complete \file{spin-noise-network} distribution: the \ts{}
acquisition kit, the uploader, the metadata schema, the protocol and policy
documents, and testing tools. Only two files go onto the spectrometer
workstation (next section):
\begin{itemize}
\item \file{topspin/spin\_noise\_run.py} --- the acquisition script (Jython,
      runs inside \ts).
\item \file{topspin/pp/zgnoise2d} --- the no-pulse pseudo-2D pulse program
      used for the noise blocks.
\end{itemize}
Keep the rest of the unpacked folder around on whatever machine will do the
upload: the uploader, \file{uploader/upload\_bundle.py}, expects
\file{schema/meta.schema.json} one level up from itself, as in the release
layout. If you move the uploader out on its own, pass
\texttt{-{}-schema <path>} explicitly --- otherwise its
\texttt{-{}-selftest} skips schema validation with a WARN yet still prints
\texttt{RESULT: PASS}.

\section{Install (two file copies)}

\begin{enumerate}
\item Copy \file{spin\_noise\_run.py} to
      \file{<TSHOME>/exp/stan/nmr/py/user/}.
      Alternative: in \ts{} type \texttt{edpy}, use \emph{File
      $\rightarrow$ Import\ldots} and pick the file --- that lands it in the
      same place.
\item Copy \file{pp/zgnoise2d} to
      \file{<TSHOME>/exp/stan/nmr/lists/pp/user/}.
\item There is no step 3. No Python packages, no licenses, no network access
      is needed on the spectrometer.
\end{enumerate}

\file{<TSHOME>} is your \ts{} installation directory ---
\file{/opt/topspin4.1.4} or \file{/opt/topspin3.6.5} on Linux,
\file{C:\textbackslash Bruker\textbackslash TopSpin4.1.4} on Windows, etc.
If you are unsure, type \texttt{set} inside \ts{} or look at the title bar of
the \ts{} window.

The script also installs \file{zgnoise2d} automatically on first run if it
can find your \ts{} directory, so step 2 is a belt-and-braces copy --- do it
anyway if you can.

\textbf{Verify \ts{} sees the script:} type \texttt{edpy} and
\file{spin\_noise\_run} should appear in the \emph{user} source directory
listing. The pulse program is checked again at run time: if \texttt{zg}
cannot find it, the script tells you and detects the file once you copy it in
by hand (Section~\ref{sec:trouble}).

\section{Verify without touching hardware}
\label{sec:verify}

Two test modes let you exercise the whole flow before the magnet is involved.
Both are started from the \ts{} command line with any \hnmr{} dataset open (a
PROTON demo set is fine --- it is only a parameter template).

\subsection*{5.1\quad SIMULATE --- the pure dry run}

\begin{lstlisting}
xpy spin_noise_run simulate
\end{lstlisting}

SIMULATE skips all spectrometer commands (and the setup blocks that lead to
them) while still running every dialog, the dataset bookkeeping, the
\file{meta.json} writing, and the zip bundling. Answer the dialogs with test
values (suggested: institution \texttt{Desk Test Lab}, slug
\texttt{desktest-lab}, sample \texttt{distilled water}, \htwo{} \texttt{100},
\dtwo{} \texttt{0}, duration \texttt{30 min}, lock \texttt{OFF}, sweep
confirmed \texttt{OFF}). Success looks like:
\begin{itemize}
\item the greeting dialog shows \texttt{*** SIMULATE MODE ... ***}.
\item all operator dialogs appear, in the order listed in
      Section~\ref{sec:dialogs}.
\item the terminal shows \texttt{SIMULATE -> ...} / \texttt{SIMULATE: zg
      mocked ...} lines and no errors.
\item the final dialog reports the bundle path, and the zip exists there.
\end{itemize}

\subsection*{5.2\quad DESKTEST --- the stronger check}

\begin{lstlisting}
xpy spin_noise_run desktest
\end{lstlisting}

DESKTEST drives the \emph{real} \ts{} API --- the full dialog chain,
\texttt{GETPAR}/\texttt{PUTPAR} on live parameter sets, \texttt{WR}/\texttt{RE}
dataset creation, \file{meta.json} writing, zip bundling --- and mocks only
the five hardware commands (\texttt{atma}, \texttt{topshim},
\texttt{pulsecal}, \texttt{zg}, \texttt{rga}). It runs fine on a
processing-only \ts{} install (Bruker's free academic license), so you can do
it at your desk. Pass criteria, from the release's desk-test checklist:
\begin{itemize}
\item terminal shows \texttt{DESKTEST -> mocked 'atma'}, \texttt{'topshim'},
      \texttt{'pulsecal'}, \texttt{'rga'} --- and no manual-fallback dialog
      appears for any of them, and no Jython traceback anywhere.
\item a dataset \file{SPINNOISE\_<yyyymmdd>\_<hhmm>} exists with expnos
      \textbf{1, 10, 14, 15, 16, 11, 12, 13}.
\item \file{meta.json} is written twice (dataset directory and bundle) and
      contains a \texttt{software} block with
      \texttt{"script\_version": "0.3.0"} (this release),
      \texttt{"schema\_version": "1.2"} (the orchestrator's Bruker-only
      metadata vintage, accepted under the repository's vendor-neutral v2.0
      schema), a \texttt{script\_sha256} fingerprint, and
      \texttt{"run\_mode": "desktest"} --- plus a \texttt{clock\_audit} block
      with one timing entry per acquisition.
\item the bundle zip,
      \file{SPINNOISE\_<date>\_<time>/}\allowbreak\file{spinnoise\_<slug>\_<YYYYMMDD\_HHMMSS>Z\_<hex>.zip},
      opens in any zip tool, and the validator passes on any Python-3 machine:
\begin{lstlisting}
python3 uploader/upload_bundle.py <bundle.zip> --selftest
      -> RESULT: PASS
\end{lstlisting}
\end{itemize}
\textbf{Never upload a SIMULATE or DESKTEST bundle} --- it contains copies of
your demo dataset, not noise data. \file{meta.json} flags it, but do not rely
on the flag --- just do not send it.

\section{The real run}
\label{sec:dialogs}

One honest note before the first live run: as of v0.3.0 the script has been
desk-tested end to end (a real Jython interpreter against a stubbed \ts{}
API, plus the SIMULATE/DESKTEST modes above), but it has not yet run on real
spectrometer hardware. We therefore treat the first live run at each facility
as a supervised pilot --- write to
\href{mailto:jwbquantum@gmail.com}{jwbquantum@gmail.com} and we will
screen-share through it.

\subsection*{6.1\quad Before you type the command}

Insert the water tube (sample changer or manual --- the script never ejects
or injects), set your usual VT setpoint, and let the sample equilibrate for a
few minutes. Tune/match and shimming happen \emph{inside} the run:
the script calls \texttt{atma} and \texttt{topshim} where your system has
them and walks you through the manual equivalents (\texttt{wobb}, your usual
shim routine) where it does not. Open any existing \hnmr{} dataset, then:
\begin{lstlisting}
xpy spin_noise_run
\end{lstlisting}

\subsection*{6.2\quad The dialogs, in order}

\begin{enumerate}
\item \dialog{Welcome} --- summarizes what the run will do (Start/Cancel).
\item \dialog{Your facility (1/5)} --- institution, city, country, contact
      e-mail, so the bundle is attributed correctly.
\item \dialog{Facility slug (2/5)} --- short machine-readable ID
      (auto-suggested), used in the bundle file name.
\item \dialog{Contact consent} --- whether the maintainers may contact you at
      that address. The e-mail is stored only if you answer yes.
\item \dialog{The sample (3/5)} --- description, \htwo{} fraction (\%),
      \dtwo{} fraction (\%), additives, tube OD, volume. Answer exactly as the
      tube is. Tap water with a described history beats an undocumented
      perfect sample.
\item \dialog{Temperature (4/5)} --- VT setpoint in K (enter 298 if VT is off
      and the bore is at room temperature).
\item \dialog{Noise-block duration (5/5)} --- 30\,min (default), 60\,min,
      180\,min, or overnight (10\,h). Setup and references add
      ${\sim}15$\,min on top. Longer runs average down the spectrum, and
      an overnight run on an idle weekend magnet is gold.
\item \dialog{Lock state} --- turn the lock OFF for the noise block if you
      can (\texttt{lock off} / BSMS LOCK key): lock rf leaks into the \hnmr{}
      channel on some systems. Either answer is accepted and recorded.
\item \dialog{BSMS field sweep --- the one that matters.} Open
      \texttt{bsmsdisp} and \emph{physically verify} SWEEP is OFF before
      answering. A sweeping field smears the spin-noise feature over several
      kHz and leaves a perfectly healthy-looking flat noise floor --- this
      silently contaminated an archival dataset in 2022, and the script
      cannot verify the sweep state on all console generations, so the
      confirmation is on you. If you cannot confirm, a follow-up dialog
      offers proceed-flagged or abort. Prefer fixing it to proceeding flagged.
\item \dialog{Hardware check} --- \ts{} version, console, probe string
      (pre-filled where the script can detect them) --- correct as needed.
\item \dialog{Probe type} --- room-temperature, nitrogen-cooled, helium
      cryoprobe, or unknown (this sets which side of the temperature-contrast
      law your point lands on).
\item \dialog{Probe temperatures (optional)} --- coil and preamp temperatures
      in K if you know them (\texttt{edte} / cryopanel). Blank is fine.
\item \dialog{90-degree pulse} --- confirm P90 ($\mu$s) and its power (dB),
      pre-filled from \texttt{pulsecal} or the current P1. Sanity-check it
      (${\sim}$7--15\,$\mu$s is typical) before confirming.
\item \emph{Noise block auto-start} --- after the references the script
      prints the row count, per-row time, and RG to the status line and
      starts the pulse-free block by itself 30\,s later. There is no
      dialog to click: \emph{the P90 confirmation above is the
      walk-away point} (the RG ladder and opening reference run
      unattended in between, ${\sim}$15\,min).
\item \dialog{Final notes} --- anything unusual during the run (construction,
      elevator, He fill, lock left on, \ldots). Blank if all was quiet.
\item \dialog{Done} --- the bundle path and the one-line upload command.
\end{enumerate}
Two conditional extras: if the script cannot find your \ts{} directory it
asks once for the user pulse-program path, and on systems missing
\texttt{atma}/\texttt{topshim}/\texttt{pulsecal}/\texttt{rga} it asks you to
perform that one step manually and press OK.

\subsection*{6.3\quad What it does, and what it will not touch}

The run creates one dataset, \file{SPINNOISE\_<date>\_<time>}, in your current data
directory:

\begin{center}
\begin{tabular}{@{}ll@{}}
\toprule
expno & role \\
\midrule
1 & setup: tune/match, shim, P90 calibration \\
10, 14, 15, 16 & RG ladder: quick $1^{\circ}$ 1Ds at RG $=$ 1, 8, 64, max \\
11 & reference (open): $1^{\circ}$ pseudo-2D, 8 rows $\times$ ${\sim}19$\,s \\
12 & \textbf{noise}: \file{zgnoise2d}, \emph{no pulse at all}, NS$=$1/row,
     RG max stable \\
13 & reference (close): same as 11 (drift check) \\
\bottomrule
\end{tabular}
\end{center}

Pulsing is limited to the standard pulse calibration and ${\sim}1^{\circ}$
tips at the calibrated observe power --- no long or repeated high-power
irradiation. The script never changes instrument configuration, never touches
your lock or sweep settings silently (it asks you and records your answer),
and reads and writes nothing outside the \file{SPINNOISE\_<date>\_<time>} dataset
plus the two installed files. Complete removal, if you ever want everything
gone:
\begin{lstlisting}
rm -rf <DATADIR>/SPINNOISE_<yyyymmdd>_<hhmm>
rm <TSHOME>/exp/stan/nmr/py/user/spin_noise_run.py
rm <TSHOME>/exp/stan/nmr/lists/pp/user/zgnoise2d
\end{lstlisting}
(Windows: delete the same three paths in Explorer. \file{<DATADIR>} is the
data directory of the template dataset you had open.)

The final dialog prints the full path of the finished bundle, which lands
inside the \file{SPINNOISE\_<date>\_<time>} dataset directory as
\file{spinnoise\_<slug>\_<timestamp>Z\_<hex>.zip}. You keep this local copy
no matter what happens during upload.

\section{Upload}

\subsection*{7.1\quad One-time setup}

On whichever machine will do the upload (Python $\geq$ 3.6, internet), go to
the \file{uploader/} folder, copy \file{config.example.json} to
\file{config.json}, and fill in the values \textbf{the coordinator sent your
facility privately}:
\begin{lstlisting}
{
  "endpoint_url": "https://...workers.dev/ingest",
  "token":        "<from the coordinator>",
  "facility_slug": "your-facility-slug"
}
\end{lstlisting}
The endpoint and token arrive by a private channel, never in e-mail chains
--- treat the token like a password. \file{config.json} is read locally and
never uploaded.

\subsection*{7.2\quad The upload}

\begin{lstlisting}
python3 uploader/upload_bundle.py /path/to/spinnoise_*.zip
\end{lstlisting}

Size never matters: bundles up to 50\,MiB go in one request, and anything
larger (overnight runs, up to 5\,GiB) automatically switches to a chunked
upload in 50\,MiB parts. Progress is checkpointed next to the zip
(\file{<bundle>.upload-state.json}), so if the network drops or the machine
reboots, \textbf{rerun the same command and it resumes where it left off}.
The uploader never deletes your zip.

Success ends with a line like:
\begin{lstlisting}
RECEIPT: <id>
\end{lstlisting}
That is the server confirming your bundle arrived intact (its SHA-256 was
verified on arrival) and was recorded. Note the id --- it is the reference
for any follow-up about this bundle.

\subsection*{7.3\quad First test run, or no endpoint yet}

If you have no \file{config.json} (or no internet path at all), the script
prints exactly what to do: the zip stays where it is, and you can e-mail or
file-transfer it to
\href{mailto:jwbquantum@gmail.com}{jwbquantum@gmail.com} (a share link is
fine if it exceeds mail limits). For a first test you can also validate
locally without any network:
\begin{lstlisting}
python3 uploader/upload_bundle.py <bundle.zip> --selftest
\end{lstlisting}

\section{Troubleshooting}
\label{sec:trouble}

\begin{itemize}
\item \textbf{``No dataset is open''} --- open any \hnmr{} dataset first (it
      is only the parameter template).
\item \textbf{Pulse program not found at \texttt{zg}} --- copy
      \file{pp/zgnoise2d} into the user pulse-program directory
      (\file{<TSHOME>/exp/stan/nmr/lists/pp/user/}) by hand and rerun ---
      the script will detect it.
\item \textbf{A \texttt{parmode} dialog appears} --- some \ts{} versions ask
      before converting a dataset to 2D. Answer yes/OK --- the dataset is
      fresh, there is nothing to lose.
\item \textbf{Old \ts{} (2.x)} --- everything is written for Jython 2.2-level
      syntax. If \file{Avance.incl} is missing, delete the \texttt{\#include}
      line in \file{zgnoise2d} (it is not used).
\item \textbf{No ATM probe / no \texttt{rga}} --- expected and handled: the
      script asks you to tune/match with \texttt{wobb} (or set RG) manually
      and confirm in a dialog.
\item \textbf{A hardware command hangs with no progress for over 5 minutes}
      --- kill the script window. Any acquisition already started finishes on
      its own and all data stays in \file{SPINNOISE\_<date>\_<time>}.
\item \textbf{A Jython error dialog appears} --- the run stopped, but started
      acquisitions finish and the data stays in \file{SPINNOISE\_<date>\_<time>}.
      Copy the error text verbatim into a support issue (below) --- a stuck
      run reported well is worth almost as much as a clean one.
\item \textbf{Facility Python 3 older than 3.6} --- run the selftest and
      upload from any other machine (copy the zip on a stick). The uploader
      is stdlib-only by design, so please do not pip-install anything on
      the spectrometer workstation.
\end{itemize}

\textbf{Where to get help:} open a facility-support issue at
\url{https://github.com/CASPEr-Night/spin-noise-network/issues/new?template=facility-support.yml}
(``don't know'' is always an acceptable answer on that form), or write to
\href{mailto:jwbquantum@gmail.com}{jwbquantum@gmail.com}.

\vspace{1.5em}
\noindent\rule{\textwidth}{0.4pt}\\[2pt]
{\small Data policy: your raw data remain yours, and contributing facilities
are co-authors on network publications their data enter --- full terms in
\file{DATA\_POLICY.md} and at
\href{https://night.blanchard-science.com}{night.blanchard-science.com}.}

\end{document}
